% Canonical Paper II appendix.
This appendix gives the sign-sensitive calculations behind
Propositions~\ref{prop:surface-frame-divergence} and
\ref{prop:surface-frame-laplacian} and
Theorem~\ref{thm:surface-adjoint-factorization}.  All calculations are local, but the
resulting identities are intrinsic.

\subsection{Connection coefficients from the frame bracket}

Let $e_1=\Xu$, $e_2=\Xv$, and write
\begin{equation}\label{eq:appendix-frame-bracket}
 [e_1,e_2]=c_ue_1+c_ve_2.
\end{equation}
Let $\omega$ be the Levi--Civita connection one-form in the convention
\begin{equation}\label{eq:appendix-connection-form-convention}
 \nabla_Ye_1=\omega(Y)e_2,
 \qquad
 \nabla_Ye_2=-\omega(Y)e_1.
\end{equation}
Torsion freeness gives
\begin{align*}
 [e_1,e_2]
 &=\nabla_{e_1}e_2-\nabla_{e_2}e_1\\
 &=-\omega(e_1)e_1-\omega(e_2)e_2.
\end{align*}
Comparison with~\eqref{eq:appendix-frame-bracket} yields
\begin{equation}\label{eq:appendix-connection-structure-coefficients}
 \omega(e_1)=-c_u,
 \qquad
 \omega(e_2)=-c_v.
\end{equation}
Therefore
\begin{align}
 \nabla_{e_1}e_1&=-c_u e_2,
 &\nabla_{e_1}e_2&=c_u e_1,
 \label{eq:appendix-connection-first-row}\\
 \nabla_{e_2}e_1&=-c_v e_2,
 &\nabla_{e_2}e_2&=c_v e_1.
 \label{eq:appendix-connection-second-row}
\end{align}
Taking traces gives
\begin{align*}
 \diver_ge_1
 &=\langle\nabla_{e_1}e_1,e_1\rangle
   +\langle\nabla_{e_2}e_1,e_2\rangle
 =-c_v,\\
 \diver_ge_2
 &=\langle\nabla_{e_1}e_2,e_1\rangle
   +\langle\nabla_{e_2}e_2,e_2\rangle
 =c_u.
\end{align*}
This proves the divergence signs in
\eqref{eq:surface-frame-divergences} without a coordinate choice.

\subsection{Integration by parts and the real energy operator}

For a real vector field $X$ and compactly supported complex-valued functions
$f,h$, apply the divergence theorem to $f\overline h X$:
\begin{align*}
 0
 &=\int_M\diver_g(f\overline h X)\,d\vol_g\\
 &=\int_M
   \bigl((Xf)\overline h+f\,X(\overline h)
         +f\overline h\,\diver_gX\bigr)\,d\vol_g.
\end{align*}
Since $X(\overline h)=\overline{Xh}$,
\begin{equation}\label{eq:appendix-vector-field-adjoint}
 \int_M(Xf)\overline h\,d\vol_g
 =\int_M f\,
  \overline{(-X-\diver_gX)h}\,d\vol_g.
\end{equation}
Thus $X^*=-X-\diver_gX$.  Inserting the two divergences above gives
\[
 e_1^*=-e_1+c_v,
 \qquad
 e_2^*=-e_2-c_u.
\]
Consequently,
\begin{align*}
 e_1^*e_1+e_2^*e_2
 &=-e_1^2-e_2^2+c_ve_1-c_ue_2\\
 &=-\bigl(e_1^2+e_2^2-c_ve_1+c_ue_2\bigr).
\end{align*}
On the other hand,
\begin{align*}
 \Delta_gf
 &=\diver_g\bigl((e_1f)e_1+(e_2f)e_2\bigr)\\
 &=e_1^2f+e_2^2f
   +(\diver_ge_1)e_1f+(\diver_ge_2)e_2f\\
 &=e_1^2f+e_2^2f-c_ve_1f+c_ue_2f.
\end{align*}
This verifies both
\eqref{eq:surface-frame-laplacian} and
\eqref{eq:surface-real-energy-factorization}.

\subsection{Raw complex products}

Put
\[
 A=e_1-ie_2=2\dAEG,
 \qquad
 B=e_1+ie_2=2\dbarAEG.
\]
Then
\begin{align}
 AB
 &=(e_1-ie_2)(e_1+ie_2)\notag\\
 &=e_1^2+e_2^2+i(e_1e_2-e_2e_1)\notag\\
 &=S_M+i(c_ue_1+c_ve_2),
 \label{eq:appendix-raw-product-ab}
\end{align}
whereas
\begin{equation}\label{eq:appendix-raw-product-ba}
 BA=S_M-i(c_ue_1+c_ve_2).
\end{equation}
These formulas use only the differential expressions; neither a volume measure
nor a domain has entered.

\subsection{Complex formal adjoints and cancellation of the drift}

With the complex $L^2$ inner product, scalar coefficients are conjugated when
an adjoint is taken.  Therefore
\begin{align*}
 \dbarAEG^*
 &=\frac12(e_1^*-ie_2^*)\\
 &=\frac12(-e_1+c_v+ie_2+ic_u)\\
 &=-\dAEG+\frac12(c_v+ic_u),
\end{align*}
and
\begin{align*}
 \dAEG^*
 &=\frac12(e_1^*+ie_2^*)\\
 &=-\dbarAEG+\frac12(c_v-ic_u).
\end{align*}
For the first quadratic product, use
~\eqref{eq:appendix-raw-product-ab}:
\begin{align*}
 4\dbarAEG^*\dbarAEG
 &=(-A+c_v+ic_u)B\\
 &=-AB+(c_v+ic_u)(e_1+ie_2)\\
 &=-S_M-i(c_ue_1+c_ve_2)\\
 &\hspace{1.5em}{}+c_ve_1+ic_ve_2+ic_ue_1-c_ue_2\\
 &=-S_M+c_ve_1-c_ue_2\\
 &=-\Delta_g.
\end{align*}
For the conjugate product, use
~\eqref{eq:appendix-raw-product-ba}:
\begin{align*}
 4\dAEG^*\dAEG
 &=(-B+c_v-ic_u)A\\
 &=-BA+(c_v-ic_u)(e_1-ie_2)\\
 &=-S_M+i(c_ue_1+c_ve_2)\\
 &\hspace{1.5em}{}+c_ve_1-ic_ve_2-ic_ue_1-c_ue_2\\
 &=-S_M+c_ve_1-c_ue_2\\
 &=-\Delta_g.
\end{align*}
This term-by-term cancellation is the reason the raw bracket twist and the
Riemannian drift must not be treated as competing definitions of a Laplacian.

\subsection{Gauge check}

Let
\[
 e_1'=\cos\vartheta\,e_1+\sin\vartheta\,e_2,
 \qquad
 e_2'=-\sin\vartheta\,e_1+\cos\vartheta\,e_2.
\]
Then
\begin{align*}
 e_1'+ie_2'
 &=(\cos\vartheta-i\sin\vartheta)e_1
   +(\sin\vartheta+i\cos\vartheta)e_2\\
 &=e^{-i\vartheta}(e_1+ie_2).
\end{align*}
If $(\eta^1,\eta^2)$ and $(\eta^{1'},\eta^{2'})$ are the corresponding
dual coframes, direct inversion of the rotation matrix gives
\[
 \eta^{1'}-i\eta^{2'}
 =e^{i\vartheta}(\eta^1-i\eta^2).
\]
Hence
\[
 \left(\frac12(e_1'+ie_2')F\right)
 (\eta^{1'}-i\eta^{2'})
 =
 \left(\frac12(e_1+ie_2)F\right)
 (\eta^1-i\eta^2),
\]
which is the explicit $U(1)$-gauge invariance of the intrinsic
$\bar\partial_MF$.

\subsection{Hyperbolic sign check}

On the basic hyperbolic AES,
\[
 e_1=-\mu y\partial_x,
 \qquad
 e_2=-\lambda y\partial_y,
 \qquad
 [e_1,e_2]=\lambda e_1.
\]
Thus $c_u=\lambda$ and $c_v=0$.  The general formulas reduce to
\begin{align*}
 e_1^*&=-e_1,
 &e_2^*&=-e_2-\lambda,\\
 \Delta_g&=e_1^2+e_2^2+\lambda e_2,
 &\dbarAEG^*&=-\dAEG+\frac{i\lambda}{2}.
\end{align*}
Since
\[
 e_1^2=\mu^2y^2\partial_x^2,
 \qquad
 e_2^2=\lambda^2y^2\partial_y^2
          +\lambda^2y\partial_y,
 \qquad
 \lambda e_2=-\lambda^2y\partial_y,
\]
the two first-order $y$-terms cancel and
\[
 \Delta_g
 =\mu^2y^2\partial_x^2+\lambda^2y^2\partial_y^2,
\]
in agreement with the coordinate calculation in Paper~I.
